% Part: lambda-calculus
% Chapter: lambda-definability
% Section: partial-recursive-functions

\documentclass[../../../include/open-logic-section]{subfiles}

\begin{document}

\olfileid{lam}{ldf}{par}
\olsection{આંશિક પુનરાવર્તી વિધેયો \usetoken{S}{lambda definable} છે}

આંશિક પુનરાવર્તી વિધેયો મૂળભૂત વિધેયોમાંથી
સંયોજન, આદિમ પુનરાવર્તન અને અનિયત મર્યાદા
સુધીની શોધક્રિયા દ્વારા મળે છે. સામાન્ય પુનરાવર્તી
વિધેયોથી તેમનો તફાવત એ છે કે અનિયત મર્યાદા સુધીની
શોધમાં વપરાતાં વિધેયો નિયમિત હોવાં જરૂરી નથી.
નિયમિતતાની શરત ન હોવાથી લઘુત્તમીકરણ દ્વારા
વ્યાખ્યાયિત વિધેય ક્યારેક અવ્યાખ્યાયિત રહી શકે છે.

પ્રથમ નજરે એવું લાગી શકે કે (સંપૂર્ણ) સામાન્ય પુનરાવર્તી
વિધેયો બધા !!{lambda definable} છે એવું બતાવતી એ જ
રીતથી બધા આંશિક પુનરાવર્તી વિધેયો પણ
!!{lambda definable} છે એવું સાબિત થઈ જશે.
દાખલા તરીકે, $f$ અને $g$ને અનુક્રમે $F$ અને~$G$
પદો !!{lambda define} કરતાં હોય, તો $f$નું $g$ સાથેનું
સંયોજન $\lambd[x][F (G x)]$ વડે !!{lambda define}
થાય છે. પરંતુ વિધેયો આંશિક હોય ત્યારે આમાં
મુશ્કેલી આવે છે. $g(x)$ અવ્યાખ્યાયિત હોય ત્યારે
$G x$નું સામાન્ય સ્વરૂપ હોતું નથી. મોટાભાગના
કિસ્સામાં $F (G x)$નું પણ સામાન્ય સ્વરૂપ હોતું
નથી, અને આપણને એવું જ જોઈએ છે. પરંતુ
$F$ એ $\lambd[x][\lambd[y][y]]$ હોય તે કિસ્સો
વિચારો: ત્યારે $F (G x)$નું સામાન્ય સ્વરૂપ
$\lambd[y][y]$ છે.

આ મુશ્કેલી દૂર કરી શકાય છે. બધા આંશિક પુનરાવર્તી
વિધેયોને એ રીતે !!{lambda define} કરી શકાય છે કે
અવ્યાખ્યાયિત મૂલ્યોનું નિરૂપણ સામાન્ય સ્વરૂપ
વિનાનાં પદો કરે. જોકે, આવી રીતો સામાન્ય પુનરાવર્તી
વિધેયો માટે આપણે લીધેલી રીત કરતાં થોડી વધુ
જટિલ અને ઓછી સહજ છે. અહીં પ્રમેય પુરાવા વિના
નોંધીએ છીએ:

\begin{thm}
  બધા આંશિક પુનરાવર્તી વિધેયો !!{lambda definable} છે.
\end{thm}

\end{document}
